%%%%%%%%%%%%%%%%%%%%%part7.tex%%%%%%%%%%%%%%%%%%%%%%%%%%%%%%%%%%
% 
% Part VII: Lifecycle, Compliance & The Future
%
%%%%%%%%%%%%%%%%%%%%%%%% Springer Nature %%%%%%%%%%%%%%%%%%%%%%

\begin{partbacktext}
\part{Lifecycle, Compliance \& The Future}
\noindent In the final part of this book, we step back from technical primitives to address the operational and regulatory reality of the Trustworthy AI Architect. We explore the full lifecycle of AI systems, the evolving global legal landscape, and the practical tools required for implementation.

In \textbf{Chapter 24}, we examine \textbf{AI Supply Chain Security}, focusing on Model Provenance and the Software Bill of Materials (SBOM). \textbf{Chapter 25} introduces the \textbf{AI Trust Threat Matrix (ATTM)}, providing a structured framework for gap analysis across the Trustworthy AI Stack. \textbf{Chapter 26} analyzes the technical mechanisms for \textbf{Trustworthy MLOps}, including continuous monitoring and drift detection. \textbf{Chapter 27} provides a comparative guide to \textbf{Global Compliance}, including the GDPR, the EU AI Act, and Vietnam's Decree 13. \textbf{Chapter 28} introduces the \textbf{Practitioner's Toolkit}, ranging from Opacus to Flower. Finally, in \textbf{Chapter 29}, we close the book with a \textbf{Roadmap} for the Quantum Era and a final call to action.

\begin{casestudy} 
The 2025 "Phantom Model" attack involved a major tech firm inadvertently using a base model that had been "pre-poisoned" at the weight level by a state-sponsored adversary months before the supply chain was audited. With no SBOM or provenance record for the open-source weights, the firm had no way to verify the model"s integrity before fine-tuning. This event catalyzed the Global Compliance and Supply Chain standards described in Part VII. 
\end{casestudy} 
\end{partbacktext}
